%% ----------------------------------------------------------------
%% Week-9 problem sheet (lectures: ProbabilisticInference,
%% matrixFactorisation, mcmc).  Every number in the answers was
%% checked in numpy (2026-10-28).
%% ----------------------------------------------------------------
\documentclass{sotonExamBoxes}    % use option [answers] to include them
%% ----------------------------------------------------------------
\usepackage{graphicx}
\usepackage{alltt}
\usepackage{bm}
\usepackage{amssymb}
\usepackage{amsfonts}
\usepackage{amsmath}
\thinmuskip=1mu
%% macros.tex -- shared macro definitions for the problem sheets.
%%
%% Every sheet inputs this instead of carrying its own copy. Input it after
%% the \usepackage lines: the definitions below use bm, amsmath and amssymb.
%%
%% Six macros were defined inconsistently across the sheets when this file
%% was made; the choices are marked CONFLICT below.




\usepackage{stmaryrd}

\renewcommand{\ttdefault}{pcr}

\newcommand{\tr}{\textsf{T}}

\newcommand{\e}[1]{{\rm e}^{#1}}

\newcommand{\E}[1]{{\rm e}^{{\displaystyle #1}}}

\newcommand{\logg}[1]{\log\left(#1\right)}

\newcommand{\Prob}[1]{\mathbb{P}\left(#1\right)}

\newcommand{\Step}{\mathop{\mathrm{O}\hspace{-15pt}\rule[8.5pt]{11pt}{1.2pt}\hspace{2pt}}}

\newcommand\diag{\mathop{\mathrm{diag}}}

\newcommand{\grad}{\bm{\nabla}}

\newcommand{\dd}{\mathrm{d}}

\newcommand{\class}{\mathcal{C}}

\newcommand{\data}{\mathcal{D}}

%% CONFLICT: gaussianProcess.tex had {U}{eur}{b}{n} (Euler bold);
%% every other sheet had this. Matrices in that one sheet now look
%% different -- see the report.

\DeclareMathAlphabet{\mat}{OT1}{cmss}{bx}{n}

%% The identity needs to be distinguishable. In cmss a bold I is a plain
%% vertical bar with no crosspieces, so \mat{I} reads as a rule rather than a
%% letter. Euler's bold I keeps its serifs, so the identity alone is set in
%% Euler while everything else stays cmss.
\DeclareMathAlphabet{\matEuler}{U}{eur}{b}{n}
\newcommand{\Id}{\matEuler{I}}

%% CONFLICT and BUG FIX: 16 sheets defined this using \qpartial,
%% which is defined nowhere, so \pd was broken in all of them.
%% Only problem2.tex had the working \partial form, kept here.

\newcommand{\pd}[2]{\ensuremath{\frac{\partial #1}{\partial #2}}\xspace}

\newcommand{\M}[1]{\ensuremath{\boldsymbol{#1}}\xspace}

\newcommand{\bst}{\rule{0mm}{5mm}}

%% CONFLICT: some sheets omitted the \! . Kept the tighter spacing,
%% which is what every sheet that actually uses \av had.

\newcommand{\av}[2][\,]{\mathbb{E}_{#1}\!\left[ {\strut #2} \right]}

\newcommand{\len}[1]{\| #1 \|}

\newcommand{\inner}[2]{\left\langle #1, #2\right\rangle}

\newcommand{\squeeze}{\setlength{\itemsep}{-2pt}}

%% CONFLICT: two spacings, identical output in maths mode.

\newcommand{\bra}[1]{\left( #1 \right)}

\newcommand{\pred}[1]{\left\llbracket { \small #1} \right\rrbracket}

\newcommand{\mln}{\hat{f}_n(\bm{x})}


\newcommand{\Tr}{\mathop{\mathrm{tr}}\,}

%% NOT SHARED: gaussianProcess.tex redefines \det to take an argument
%% (\det{\mat{V}}), while calculus.tex and vectorSpaces.tex use the standard
%% no-argument operator (\det(\mat{X})). The two cannot coexist, so that
%% redefinition stays local to gaussianProcess.tex.

%% NOT SHARED: \ml and \mm mean different things in different sheets.
%% amgm and convexity use a no-argument \ml and \mm = \hat{m}_c(x);
%% ensembleLearning uses a one-argument \ml and \mm = \hat{m}(x).
%% Both forms are in use, so each sheet keeps its own.

\graphicspath{{figures/}}
\usepackage{parskip}
\AddToHook{env/samepage/begin}{%
\AddToHook{para/before}{\nopagebreak}%
\AddToHook{para/after}{\nopagebreak}%
}
\begin{document}
\unitTitle{Latent Variables, Matrix Factorisation and MCMC Problem Sheet}
\authors{Adam Pr\"ugel-Bennett}
\unitcode{AICE3001/AICE6002}
\semester{2}
\year{2022}
\duration{45}{0}

\maketitle

This sheet covers the lectures on probabilistic inference, matrix
factorisation and MCMC.  Question~1 practises generative and
hierarchical models and the EM algorithm, Question~2 low-rank approximation and matrix
completion, Question~3 drawing random samples, and Question~4 Markov
chains and the Metropolis algorithm.  Each ends with a short Python
part (\texttt{import numpy as np} is assumed).  Question~5 is harder
and optional: it proves that EM and MCMC work, and runs alternating
least squares.

% ***************************** question 1 **************************

\begin{question}{20}
  \begin{qparts}

    \qpart[3]{An email is spam with probability $0.4$.  The word
      ``free'' appears in 30\% of spam emails and in 2\% of other
      emails.  Write down the joint distribution of $S$ (spam or not)
      and $F$ (contains ``free'' or not) as a $2\times2$ table.  Find
      $\Prob{F=\text{yes}}$ and $\Prob{S=\text{spam}\mid F=\text{yes}}$.
      \begin{answer}
        Using $\Prob{S,F} = \Prob{F\mid S}\,\Prob{S}$:
        \begin{center}
          \begin{tabular}{l|cc}
            & $F=$ yes & $F=$ no \\ \hline
            $S=$ spam & $0.4\times0.3 = 0.12$ & $0.4\times0.7=0.28$ \\
            $S=$ not spam & $0.6\times0.02 = 0.012$ & $0.6\times0.98 = 0.588$
          \end{tabular}
        \end{center}
        Marginalising over $S$, $\Prob{F=\text{yes}} = 0.12+0.012 =
        0.132$, so $\Prob{S=\text{spam}\mid F=\text{yes}} = 0.12/0.132
        \approx 0.91$.  We modelled how emails are generated (a
        generative model) and used it to classify (discriminate).
      \end{answer}}{\lines{7}}

    \qpart[4]{In the hierarchical ratings model of the lecture,
      $\theta_j\sim\mathrm{Beta}(a,b)$ and $k_j\sim\mathrm{Bin}(n_j,\theta_j)$,
      the posterior mean of $\theta_j$ is $(a+k_j)/(a+b+n_j)$.  Take
      $a=20$, $b=5$.  Product $C$ has 9 positive reviews out of 10, and
      product $D$ has 85 out of 100.  Rank them by the fraction of
      positive reviews and by the posterior mean.  How many of its 10
      reviews would $C$ need to be positive to rank above $D$?
      \begin{answer}
        By fraction $C$ ($0.90$) beats $D$ ($0.85$).  By posterior mean
        $C$ gets $29/35\approx0.829$ and $D$ gets $105/125 = 0.840$, so $D$
        ranks higher.  $C$ needs $(20+k)/35 > 0.84$, i.e.\ $k > 9.4$: all
        10 reviews positive (giving $30/35\approx0.857$).  The prior is
        worth $a+b=25$ reviews at the typical rate $0.8$, which outweighs
        $C$'s ten reviews but not $D$'s hundred.
      \end{answer}}{\lines{6}}

    \qpart[3]{A latent variable $Z$ is $1$ or $0$ with probability
      $\frac12$ each.  Given $Z=1$, $X\sim\mathcal{N}(0,1)$; given $Z=0$,
      $X\sim\mathcal{N}(4,1)$.  Write down the density $f(x)$ of $X$.
      Show that the \emph{responsibility} $\gamma(x) = \Prob{Z=1\mid X=x}$
      is $1/(1+\e{4x-8})$, and evaluate it at $x=1$, $2$ and $3$.
      \begin{answer}
        Marginalising over $Z$, $f(x) = \frac12\mathcal{N}(x\mid0,1) +
        \frac12\mathcal{N}(x\mid4,1)$.  By Bayes' rule (the $\frac12$s and
        the normalisations $1/\sqrt{2\pi}$ cancel)
        \begin{align*}
          \gamma(x) = \frac{\e{-x^2/2}}{\e{-x^2/2} + \e{-(x-4)^2/2}}
          = \frac{1}{1 + \e{(x^2 - (x-4)^2)/2}} = \frac{1}{1+\e{4x-8}} .
        \end{align*}
        $\gamma(1)\approx0.982$, $\gamma(2) = 0.5$ and $\gamma(3)\approx
        0.018$.  This is the logistic sigmoid of a linear function of
        $x$: a generative model with equal variances leads to logistic
        regression.
      \end{answer}}{\lines{8}}

    \qpart[4]{Take the model of part (c) but treat $p=\Prob{Z=1}$,
      $\mu_A$ and $\mu_B$ as unknown (the variances stay 1).  The data
      are $X = (-1, 0, 1, 2, 4, 5)$.  Starting from $p=\frac12$,
      $\mu_A=0$, $\mu_B=4$, carry out one step of EM: compute the
      responsibilities $\gamma_i$ (use part (c)), then the new values
      \begin{align*}
        p = \frac1m\sum_{i=1}^m\gamma_i, \qquad
        \mu_A = \frac{\sum_i \gamma_i X_i}{\sum_i\gamma_i}, \qquad
        \mu_B = \frac{\sum_i (1-\gamma_i) X_i}{\sum_i(1-\gamma_i)} .
      \end{align*}
      \begin{answer}
        From part (c), $\gamma \approx (1.000, 1.000, 0.982, 0.500,
        0.000, 0.000)$ (more precisely $0.99966$ for $X=0$ and
        $0.00034$ for $X=4$).  Then $\sum_i\gamma_i \approx 3.482$, so
        $p\approx 3.482/6 \approx 0.58$, and
        \begin{align*}
          \mu_A &\approx \frac{-1 + 0 + 0.982 + 1.0 + 0.001}{3.482}
          \approx 0.28, &
          \mu_B &\approx \frac{0.018 + 1.0 + 4.0 + 5}{2.518} \approx 3.98 .
        \end{align*}
        The point $X=2$ sits midway, so half of it is given to each
        cluster: EM makes \emph{soft} assignments.
      \end{answer}}{\lines{8}}

    \qpart[2]{In the M-step the terms of $Q$ that involve $p$ are
      $\sum_{i=1}^m \bigl(\gamma_i\log(p) + (1-\gamma_i)\log(1-p)\bigr)$.
      Show that they are maximised by $p = \frac1m\sum_i\gamma_i$.
      \begin{answer}
        Setting the derivative to zero,
        $\frac1p\sum_i\gamma_i - \frac1{1-p}\sum_i(1-\gamma_i) = 0$, so
        $(1-p)\sum_i\gamma_i = p\,(m - \sum_i\gamma_i)$, giving
        $p = \frac1m\sum_i\gamma_i$: the average responsibility.  (The
        second derivative is negative, so this is a maximum.)
      \end{answer}}{\lines{5}}

    \qpart[4]{In Python, generate $m=200$ points from the model with
      $p=0.3$, $\mu_A=0$, $\mu_B=4$, start EM from $p=0.5$, $\mu_A=-1$,
      $\mu_B=1$, and print the log-likelihood
      $\sum_i\log\bigl(p\,\mathcal{N}(X_i\mid\mu_A,1) +
      (1-p)\,\mathcal{N}(X_i\mid\mu_B,1)\bigr)$ at each iteration.  Does
      it ever decrease?
      \begin{answer}
\begin{alltt}
rng = np.random.default\_rng(0)\\
m = 200\\
Z = rng.random(m) < 0.3\\
X = np.where(Z, rng.normal(0, 1, m), rng.normal(4, 1, m))\\
def N(x, mu):\\
\hspace*{2em}return np.exp(-(x - mu)**2 / 2) / np.sqrt(2*np.pi)\\
p, muA, muB = 0.5, -1.0, 1.0\\
for t in range(20):\\
\hspace*{2em}a, b = p * N(X, muA), (1 - p) * N(X, muB)\\
\hspace*{2em}print(t, np.sum(np.log(a + b)))\\
\hspace*{2em}g = a / (a + b)\hspace*{2em}\# E-step: responsibilities\\
\hspace*{2em}p = g.mean()\hspace*{2em}\# M-step\\
\hspace*{2em}muA = g @ X / g.sum()\\
\hspace*{2em}muB = (1 - g) @ X / (1 - g).sum()\\
print(p, muA, muB)
\end{alltt}
        The log-likelihood goes $-1056.4$, $-407.9$, $-379.2$,
        $-377.46$, $-377.37$, \ldots, and has settled at $-377.363$ by
        iteration 9.  It never decreases (once converged it wobbles by
        about $10^{-13}$, which is rounding error).  The estimates are
        $p=0.274$, $\mu_A=0.010$ and $\mu_B=4.050$: close to the truth,
        and closer still to this sample, in which 26.5\% of the points
        are of type $A$.
      \end{answer}}{\lines{8}}

  \end{qparts}
\end{question}
\newpage
% ***************************** question 2 **************************

\begin{question}{20}
  \begin{qparts}

    \qpart[3]{A film site has $1000$ users and $500$ films.  How many
      numbers does a rank $K=20$ factorisation $\mat{A}\mat{B}$ of the
      ratings matrix have, compared with the full matrix?  LoRA adds a
      rank $K=4$ update to a $1024\times1024$ weight matrix: what
      fraction of the weights does it train?
      \begin{answer}
        $(1000+500)\times20 = 30\,000$ against $1000\times500 =
        500\,000$, so 6\%.  For LoRA, $(1024+1024)\times4 = 8192$ against
        $1024^2 = 1\,048\,576$: about $0.8\%$.
      \end{answer}}{\lines{4}}

    \qpart[4]{For the ratings in the lecture
      \begin{center}
        \begin{tabular}{|l||c|c|c|c|c|}\hline
          Name & Titanic & Dr No & The Matrix & Spartacus & Ten \\ \hline
          Alice & 2 & \O & 4  & 5 & 2  \\ \hline
          Bob   & 5 & 3  & \O & 3 & \O \\ \hline
          Chas  & \O & 2 & 3 & 5 & 3 \\ \hline
        \end{tabular}
      \end{center}
      use $R_{ui}\approx\bar{R}_u + \bar{R}_i - \bar{R}$ to predict
      Bob's rating of The Matrix and Chas's rating of Titanic.
      \begin{answer}
        The overall mean of the 11 ratings is $\bar{R} = 37/11 \approx
        3.36$.  Bob: $\bar{R}_u = 11/3\approx3.67$ and The Matrix has
        $\bar{R}_i = 3.5$, giving $3.67+3.5-3.36\approx3.80$.  Chas:
        $\bar{R}_u = 13/4 = 3.25$ and Titanic has $\bar{R}_i = 3.5$,
        giving $\approx3.39$.  (Bob rates generously, Chas less so.)
      \end{answer}}{\lines{5}}

    \qpart[4]{The matrix $\mat{X} = \begin{pmatrix} 3 & 1 \\ 1 & 3
      \end{pmatrix}$ is symmetric with positive eigenvalues, so its SVD
      is its eigen-decomposition.  Find it, write down the best rank-1
      approximation $\mat{X}_1$, and check that
      $\|\mat{X}-\mat{X}_1\|_F$ equals the discarded singular value, as
      Eckart--Young says.
      \begin{answer}
        The eigenvectors are $\bm{u}_1 = (1,1)^\tr/\sqrt2$ with
        eigenvalue $s_1=4$ and $\bm{u}_2 = (1,-1)^\tr/\sqrt2$ with
        $s_2=2$.  Then
        \begin{align*}
          \mat{X}_1 = s_1\bm{u}_1\bm{u}_1^\tr = \begin{pmatrix} 2 & 2 \\
            2 & 2 \end{pmatrix}, \qquad
          \mat{X} - \mat{X}_1 = \begin{pmatrix} 1 & -1 \\ -1 & 1
          \end{pmatrix},
        \end{align*}
        and $\|\mat{X}-\mat{X}_1\|_F = \sqrt{1+1+1+1} = 2 = s_2$.  (The
        spectral norm of the error is also $2$.)
      \end{answer}}{\lines{7}}

    \qpart[4]{In alternating least squares with $K=1$ every factor is
      a number.  A user has rated three films whose factors are
      $b_j = 1$, $2$ and $0.5$, giving ratings $X_j = 2$, $4$ and $1$.
      With the item factors fixed, the user's factor is the ridge
      regression solution $a = \sum_j b_jX_j / \bigl(\eta + \sum_j
      b_j^2\bigr)$.  Find $a$ and the predicted rating of a fourth film
      with $b=1.5$ for $\eta = 0$, $0.5$ and $5$.  Comment.
      \begin{answer}
        $\sum_j b_jX_j = 2+8+0.5 = 10.5$ and $\sum_j b_j^2 = 5.25$.
        \begin{center}
          \begin{tabular}{c|ccc}
            $\eta$ & 0 & 0.5 & 5 \\ \hline
            $a$ & 2.000 & 1.826 & 1.024 \\
            prediction $1.5a$ & 3.00 & 2.74 & 1.54
          \end{tabular}
        \end{center}
        With $\eta=0$ the ratings are fitted exactly ($X_j = 2b_j$).  The
        regulariser (the prior) shrinks $a$ towards 0, more strongly
        when there are few ratings.  Shrinking towards a rating of 0 is
        a poor default, which is why in practice we first subtract
        means such as those of part (b) and factorise what is left.
      \end{answer}}{\lines{7}}

    \qpart[5]{In Python, make an $8\times6$ matrix of rank 2, hide about
      a quarter of its entries, and fill them in with the iterative SVD
      method of the lecture (start from the user means; keep the best
      rank-2 approximation; copy its values into the missing entries
      only).  Print the RMS error on the missing entries as you go.
      \begin{answer}
\begin{alltt}
rng = np.random.default\_rng(0)\\
F = rng.integers(1, 4, (8, 2))\\
X = F @ rng.integers(0, 3, (2, 6))\hspace*{2em}\# rank 2\\
W = rng.random(X.shape) > 0.25\hspace*{2em}\# True where observed\\
miss = np.logical\_not(W)\\
Xt = np.where(W, X, 0.0)\\
mean = Xt.sum(1) / W.sum(1)\hspace*{2em}\# user means\\
Xt = np.where(W, X, mean[:, None])\\
for t in range(201):\\
\hspace*{2em}if t in (0, 1, 5, 20, 50, 100, 200):\\
\hspace*{4em}print(t, np.sqrt(np.mean((Xt - X)[miss]**2)))\\
\hspace*{2em}U, s, Vt = np.linalg.svd(Xt)\\
\hspace*{2em}L = (U[:, :2] * s[:2]) @ Vt[:2]\hspace*{2em}\# rank 2\\
\hspace*{2em}Xt = np.where(W, X, L)
\end{alltt}
        Here 12 of the 48 entries are missing.  The RMS error is $2.79$
        with the user means, then $2.04$, $1.26$, $0.67$, $0.14$,
        $0.011$ and $8\times10^{-5}$ after 1, 5, 20, 50, 100 and 200
        iterations.  Because the matrix really has rank 2 the missing
        entries are recovered exactly, but convergence is slow.
      \end{answer}}{\lines{8}}

  \end{qparts}
\end{question}
\newpage
% ***************************** question 3 **************************

\begin{question}{20}
  \begin{qparts}

    \qpart[4]{We estimate $I = \int_0^1\e{x}\,\dd x = \av{\e{U}}$, with
      $U$ uniform on $(0,1)$, by averaging $\e{U_i}$ over $n$
      independent samples.  Find $\mathrm{Var}(\e{U})$, the standard
      error for $n=10^4$, and the $n$ needed for a standard error of
      $10^{-4}$.
      \begin{answer}
        $I = \e{}-1\approx1.718$ and $\av{\e{2U}} = (\e{2}-1)/2 \approx
        3.195$, so $\mathrm{Var}(\e{U}) = 3.195 - 1.718^2 \approx 0.242$.
        The standard error is $\sqrt{0.242/n}$: $0.0049$ for $n=10^4$.
        For $10^{-4}$ we need $n = 0.242/10^{-8}\approx2.4\times10^7$.
        The error falls only as $1/\sqrt{n}$, so each extra decimal place
        costs 100 times as many samples.
      \end{answer}}{\lines{6}}

    \qpart[4]{The Cauchy density is $f(x) = 1/\bigl(\pi(1+x^2)\bigr)$.
      Find its cumulative distribution function and hence a
      transformation that turns a uniform deviate $U$ into a Cauchy
      deviate.  Why would averaging Cauchy samples be a bad way to
      estimate its mean?
      \begin{answer}
        $F(x) = \int_{-\infty}^x \frac{\dd t}{\pi(1+t^2)} = \frac12 +
        \frac1\pi\arctan(x)$.  Solving $U = F(X)$ gives $X =
        \tan\bigl(\pi(U-\frac12)\bigr)$.  The Cauchy distribution has
        such thick tails that its mean and variance do not exist: the
        sample average never settles down (it is itself Cauchy
        distributed, however many samples we take).  The $1/\sqrt{n}$
        rule needs a finite variance.
      \end{answer}}{\lines{6}}

    \qpart[5]{We want samples from $f(x) = 6x(1-x)$ on $[0,1]$ using
      uniform proposals, $g(y) = 1$.  Find the smallest $c$ with
      $c\,g(x)\geq f(x)$, describe the rejection method, and find the
      fraction of proposals that are accepted.  Explain why the
      accepted values have density $f$.
      \begin{answer}
        $f$ is largest at $x=\frac12$, where $f=\frac32$, so $c=1.5$.
        Draw $Y$ and $U$ uniformly on $(0,1)$; accept $Y$ if $U\,c\,g(Y)
        \leq f(Y)$, i.e.\ with probability $f(Y)/(c\,g(Y)) = 4Y(1-Y)$;
        otherwise try again.  The density of proposing $y$ and then
        accepting it is $g(y)\,f(y)/(c\,g(y)) = f(y)/c$, which is
        proportional to $f$.  Integrating, a proposal is accepted with
        probability $1/c = \frac23$.
      \end{answer}}{\lines{7}}

    \qpart[3]{To draw samples from $\mathcal{N}(\bm{0},\mat{I})$ in $d$
      dimensions we use proposals from the wider
      $\mathcal{N}(\bm{0},s^2\mat{I})$ with $s>1$.  Show that the
      smallest $c$ is $s^d$.  What fraction of proposals is accepted for
      $s=1.2$ when $d=1$, $10$ and $100$?
      \begin{answer}
        The ratio of the densities is
        \begin{align*}
          \frac{f(\bm{x})}{g(\bm{x})} = s^d\,\e{-\frac12\len{\bm{x}}^2(1-1/s^2)} \leq s^d,
        \end{align*}
        with equality at $\bm{x}=\bm{0}$, so $c=s^d$ and the acceptance
        rate is $s^{-d}$: $0.83$, $0.16$ and $1.2\times10^{-8}$.  Even a
        slightly wrong proposal is hopeless in high dimensions, which is
        why we need MCMC.
      \end{answer}}{\lines{6}}

    \qpart[4]{In Python, implement the rejection sampler of part (c)
      with $10^5$ proposals.  Check the acceptance rate, and that the
      accepted samples have mean $\frac12$ and variance $\frac1{20}$.
      \begin{answer}
\begin{alltt}
rng = np.random.default\_rng(0)\\
n = 100000\\
Y = rng.random(n)\hspace*{2em}\# proposals from g = U(0,1)\\
U = rng.random(n)\\
x = Y[1.5 * U <= 6 * Y * (1 - Y)]\\
print(len(x) / n, x.mean(), x.var())
\end{alltt}
        This gives an acceptance rate of $0.665$ (against $\frac23$),
        mean $0.4998$ and variance $0.0496$.  Working on whole arrays at
        once is much faster in numpy than a loop that tries again after
        each rejection.
      \end{answer}}{\lines{6}}

  \end{qparts}
\end{question}
\newpage
% ***************************** question 4 **************************

\begin{question}{20}
  \begin{qparts}

    \qpart[4]{$M_{ij}$ is the probability that a Markov chain moves
      from state $j$ to state $i$.  Explain why $\sum_i M_{ij}=1$.  Show
      that if $\bm{\pi}$ satisfies detailed balance, $M_{ij}\pi_j =
      M_{ji}\pi_i$, then $\mat{M}\bm{\pi}=\bm{\pi}$.  What does this
      mean?
      \begin{answer}
        From state $j$ the chain must go somewhere (possibly staying at
        $j$), so the probabilities of all destinations sum to 1.
        Summing detailed balance over $j$,
        \begin{align*}
          (\mat{M}\bm{\pi})_i = \sum_j M_{ij}\pi_j = \sum_j M_{ji}\pi_i
          = \pi_i \sum_j M_{ji} = \pi_i .
        \end{align*}
        If the chain's states are distributed according to $\bm{\pi}$,
        they still are after another step: $\bm{\pi}$ is a stationary
        distribution (an eigenvector of $\mat{M}$ with eigenvalue 1).
      \end{answer}}{\lines{7}}

    \qpart[6]{Four states sit on a ring, $1$--$2$--$3$--$4$--$1$, with
      target probabilities $\bm{\pi} = (0.1, 0.2, 0.3, 0.4)$.  From each
      state the Metropolis algorithm proposes one of its two neighbours,
      each with probability $\frac12$, and accepts with probability
      $\min(1, \pi_i/\pi_j)$.  Write down $\mat{M}$ and check that
      $\mat{M}\bm{\pi}=\bm{\pi}$.
      \begin{answer}
        From 1: both neighbours (2 and 4) are more probable, so it moves
        to each with probability $\frac12$.  From 2: to 3 with
        $\frac12$, to 1 with $\frac12\times\frac{0.1}{0.2} = \frac14$,
        stays with $\frac14$.  From 3: to 4 with $\frac12$, to 2 with
        $\frac12\times\frac23 = \frac13$, stays with $\frac16$.  From 4:
        to 3 with $\frac12\times\frac34 = \frac38$, to 1 with
        $\frac12\times\frac14 = \frac18$, stays with $\frac12$.  Each
        column is a ``from'' state:
        \begin{align*}
          \mat{M} = \begin{pmatrix}
            0 & \frac14 & 0 & \frac18 \\[2pt]
            \frac12 & \frac14 & \frac13 & 0 \\[2pt]
            0 & \frac12 & \frac16 & \frac38 \\[2pt]
            \frac12 & 0 & \frac12 & \frac12
          \end{pmatrix},
          \qquad
          \mat{M}\bm{\pi} = \begin{pmatrix}
            0.05 + 0.05 \\ 0.05 + 0.05 + 0.1 \\ 0.1 + 0.05 + 0.15 \\
            0.05 + 0.15 + 0.2
          \end{pmatrix} = \begin{pmatrix} 0.1 \\ 0.2 \\ 0.3 \\ 0.4
          \end{pmatrix}.
        \end{align*}
      \end{answer}}{\lines{10}}

    \qpart[3]{In Python, find the eigenvalues of $\mat{M}$, and starting
      from state 1, $\bm{p}(0) = (1,0,0,0)^\tr$, print
      $\bm{p}(t) = \mat{M}^t\bm{p}(0)$ for $t$ up to 10.  How quickly does
      it approach $\bm{\pi}$?
      \begin{answer}
\begin{alltt}
M = np.array([[0, 1/4, 0, 1/8], [1/2, 1/4, 1/3, 0],\\
\hspace*{6em}[0, 1/2, 1/6, 3/8], [1/2, 0, 1/2, 1/2]])\\
print(np.linalg.eigvals(M))\\
p = np.array([1.0, 0, 0, 0])\\
for t in range(11):\\
\hspace*{2em}print(t, p)\\
\hspace*{2em}p = M @ p
\end{alltt}
        The eigenvalues are $1$, $0.386$, $0.052$ and $-0.521$.  After
        10 steps $\bm{p} = (0.100, 0.199, 0.301, 0.399)$.  The error is
        dominated by the eigenvalue $-0.521$: it roughly halves each step
        and changes sign, so $\bm{p}(t)$ overshoots back and forth as it
        converges.
      \end{answer}}{\lines{6}}

    \qpart[3]{To sample from a posterior $\pi(\bm{\theta}) =
      f(\data|\bm{\theta})\,f(\bm{\theta})/f(\data)$ the Metropolis
      algorithm accepts a proposal $\bm{\theta}'$ with probability
      $\min(1, \pi(\bm{\theta}')/\pi(\bm{\theta}))$.  Why don't we need
      $f(\data)$?  Why do programs compute $\log\pi(\bm{\theta}') -
      \log\pi(\bm{\theta})$ rather than the ratio itself?
      \begin{answer}
        $f(\data)$ appears in both $\pi(\bm{\theta}')$ and
        $\pi(\bm{\theta})$ and cancels from the ratio, so we only need
        likelihood times prior.  The likelihood is a product of many
        probabilities and soon underflows to 0 in floating point; its
        log is a sum and causes no trouble.  (Accept if $\log U <
        \log\pi(\bm{\theta}') - \log\pi(\bm{\theta})$.)
      \end{answer}}{\lines{5}}

    \qpart[4]{In Python, use Metropolis with proposals $x' = x +
      s\,\epsilon$, $\epsilon\sim\mathcal{N}(0,1)$, to sample from
      $\pi(x)\propto\e{-(x^2-4)^2/8}$, which has peaks at $x=\pm2$.  Start
      at $x=2$ and run $10^5$ steps for $s = 0.1$, $1$ and $10$.  For each,
      print the acceptance rate, the fraction of samples with $x>0$
      (which should be $\frac12$), the number of times $x$ changes sign,
      and the estimate of $\av{x^2}$ (whose exact value is $3.41$).
      Comment.
      \begin{answer}
\begin{alltt}
rng = np.random.default\_rng(0)\\
log\_pi = lambda x: -(x**2 - 4)**2 / 8\\
for s in [0.1, 1.0, 10.0]:\\
\hspace*{2em}x, xs, acc = 2.0, [], 0\\
\hspace*{2em}for t in range(100000):\\
\hspace*{4em}y = x + s * rng.standard\_normal()\\
\hspace*{4em}if np.log(rng.random()) < log\_pi(y) - log\_pi(x):\\
\hspace*{6em}x, acc = y, acc + 1\\
\hspace*{4em}xs.append(x)\hspace*{2em}\# repeats count when we stay\\
\hspace*{2em}xs = np.array(xs)\\
\hspace*{2em}sg = np.sign(xs)\\
\hspace*{2em}crossings = np.sum(sg[1:] != sg[:-1])\\
\hspace*{2em}print(s, acc / len(xs), np.mean(xs > 0))\\
\hspace*{2em}print(crossings, np.mean(xs**2))
\end{alltt}
        \begin{center}
          \begin{tabular}{c|cccc}
            $s$ & accepted & $x>0$ & sign changes & $\av{x^2}$ \\ \hline
            0.1 & 0.95 & 0.61 & 221 & 3.47 \\
            1 & 0.59 & 0.51 & 4533 & 3.41 \\
            10 & 0.15 & 0.49 & 7257 & 3.42
          \end{tabular}
        \end{center}
        With $s=0.1$ almost every move is accepted, but the chain creeps
        and rarely crosses the dip at $x=0$, so it spends too long in the
        peak where it started.  With $s=10$ most proposals are rejected,
        but those accepted often jump straight between the peaks.  A step
        size comparable with the width of $\pi$ works well.
      \end{answer}}{\lines{8}}

  \end{qparts}
\end{question}
\newpage

% ***************************** question 5 **************************

\begin{question}{25}
  This question is harder and optional.  It shows why EM never
  decreases the likelihood, why an MCMC chain with detailed balance
  settles down, and runs alternating least squares on the lecture's
  ratings.
  \begin{qparts}

    \qpart[7]{For data $\data$, latent variables $\bm{Z}$ and parameters
      $\bm{\Theta}$, show that for every $\bm{Z}$
      \begin{align*}
        \log f(\data|\bm{\Theta}) = \log f(\data,\bm{Z}|\bm{\Theta}) -
        \log\Prob{\bm{Z}|\data,\bm{\Theta}} .
      \end{align*}
      Average both sides over $P_t(\bm{Z}) =
      \Prob{\bm{Z}|\data,\bm{\Theta}^{(t)}}$ to get $\log
      f(\data|\bm{\Theta}) = Q(\bm{\Theta}|\bm{\Theta}^{(t)}) +
      S(\bm{\Theta})$, where $Q$ is the function EM maximises and $S(\bm{\Theta}) =
      -\sum_{\bm{Z}}P_t(\bm{Z})\log\Prob{\bm{Z}|\data,\bm{\Theta}}$.  Using
      $\log(x)\leq x-1$, show that $S(\bm{\Theta})\geq S(\bm{\Theta}^{(t)})$
      for every $\bm{\Theta}$, and hence that an EM step cannot decrease
      the likelihood.
      \begin{answer}
        The identity is $f(\data,\bm{Z}|\bm{\Theta}) =
        \Prob{\bm{Z}|\data,\bm{\Theta}}\,f(\data|\bm{\Theta})$ (the
        definition of conditional probability) after taking logs.  The
        left-hand side does not depend on $\bm{Z}$ and
        $\sum_{\bm{Z}}P_t(\bm{Z}) = 1$, so averaging gives
        \begin{align*}
          \log f(\data|\bm{\Theta}) = \underbrace{\sum_{\bm{Z}}P_t(\bm{Z})
          \log f(\data,\bm{Z}|\bm{\Theta})}_{Q(\bm{\Theta}|\bm{\Theta}^{(t)})}
          \underbrace{-\sum_{\bm{Z}}P_t(\bm{Z})\log
          \Prob{\bm{Z}|\data,\bm{\Theta}}}_{S(\bm{\Theta})} .
        \end{align*}
        Writing $P(\bm{Z}) = \Prob{\bm{Z}|\data,\bm{\Theta}}$,
        \begin{align*}
          S(\bm{\Theta}^{(t)}) - S(\bm{\Theta}) = \sum_{\bm{Z}} P_t(\bm{Z})
          \log\frac{P(\bm{Z})}{P_t(\bm{Z})} \leq \sum_{\bm{Z}} P_t(\bm{Z})
          \left(\frac{P(\bm{Z})}{P_t(\bm{Z})} - 1\right) = 1 - 1 = 0 .
        \end{align*}
        (The difference $S(\bm{\Theta})-S(\bm{\Theta}^{(t)})$ is the KL
        divergence between $P_t$ and $P$.)  The M-step chooses
        $\bm{\Theta}^{(t+1)}$ to maximise $Q$, so
        $Q(\bm{\Theta}^{(t+1)}|\bm{\Theta}^{(t)})\geq
        Q(\bm{\Theta}^{(t)}|\bm{\Theta}^{(t)})$, and $S$ cannot go down
        either, so $\log f(\data|\bm{\Theta}^{(t+1)}) \geq \log
        f(\data|\bm{\Theta}^{(t)})$.
      \end{answer}}{\lines{14}}

    \qpart[6]{Let $\mat{D} = \diag(\bm{\pi})$ with all $\pi_i>0$.  Show
      that if $\mat{M}$ satisfies detailed balance then
      $\mat{A} = \mat{D}^{-1/2}\mat{M}\mat{D}^{1/2}$ is symmetric.
      Deduce that the eigenvalues of $\mat{M}$ are real and that
      $\mat{M}$ has a full set of eigenvectors (the assumption made in
      the lecture when we expanded $\bm{p}(0)$ in eigenvectors).
      Check that $A_{12} = A_{21}$ for the matrix of Question~4(b).
      \begin{answer}
        $A_{ij} = \pi_i^{-1/2}M_{ij}\pi_j^{1/2}$ and $A_{ji} =
        \pi_j^{-1/2}M_{ji}\pi_i^{1/2}$.  Multiplying both by
        $\sqrt{\pi_i\pi_j}$ gives $M_{ij}\pi_j$ and $M_{ji}\pi_i$, which
        are equal by detailed balance, so $A_{ij} = A_{ji}$.  A real
        symmetric matrix has real eigenvalues and an orthonormal set of
        eigenvectors $\bm{w}_k$.  If $\mat{A}\bm{w} = \lambda\bm{w}$ then
        $\mat{M}(\mat{D}^{1/2}\bm{w}) = \mat{D}^{1/2}\mat{A}\bm{w} =
        \lambda\,\mat{D}^{1/2}\bm{w}$, so $\mat{M}$ has the same
        eigenvalues, with independent eigenvectors $\mat{D}^{1/2}\bm{w}_k$.
        For Question~4(b), $A_{12} = \frac14\sqrt{0.2/0.1} =
        \frac{\sqrt2}{4}$ and $A_{21} = \frac12\sqrt{0.1/0.2} =
        \frac{1}{2\sqrt2} = \frac{\sqrt2}{4}$.
      \end{answer}}{\lines{10}}

    \qpart[5]{In matrix factorisation with $\mat{B}$ fixed, row
      $\bm{a}_i$ of $\mat{A}$ minimises $\sum_j (X_{ij} -
      \bm{a}_i^\tr\bm{b}_j)^2 + \eta\len{\bm{a}_i}^2$, the sum running
      over the films user $i$ has rated.  Differentiate to show that
      $\bm{a}_i = \bigl(\eta\,\mat{I} + \sum_j\bm{b}_j\bm{b}_j^\tr
      \bigr)^{-1}\sum_j\bm{b}_jX_{ij}$, and check that it reduces to the
      formula of Question~2(d) when $K=1$.  Why is the matrix being
      inverted always invertible when $\eta>0$?
      \begin{answer}
        The gradient is $-2\sum_j\bm{b}_j(X_{ij} - \bm{b}_j^\tr\bm{a}_i)
        + 2\eta\,\bm{a}_i$ (using $\bm{a}_i^\tr\bm{b}_j =
        \bm{b}_j^\tr\bm{a}_i$).  Setting it to zero gives
        $\bigl(\eta\,\mat{I} + \sum_j\bm{b}_j\bm{b}_j^\tr\bigr)\bm{a}_i =
        \sum_j\bm{b}_jX_{ij}$.  For $K=1$ this is $(\eta + \sum_jb_j^2)\,a
        = \sum_jb_jX_j$.  For any $\bm{v}\neq\bm{0}$,
        $\bm{v}^\tr\bigl(\eta\,\mat{I} + \sum_j\bm{b}_j\bm{b}_j^\tr\bigr)\bm{v}
        = \eta\len{\bm{v}}^2 + \sum_j(\bm{b}_j^\tr\bm{v})^2 > 0$, so the
        matrix is positive definite and hence invertible, even if the user
        has rated fewer than $K$ films.
      \end{answer}}{\lines{9}}

    \qpart[7]{In Python, run alternating least squares with $K=1$ and
      $\eta=0.1$ on the ratings of Question~2(b) (use \texttt{np.nan}
      for the missing entries), starting from random factors.  Compare
      the predictions for Bob's rating of The Matrix and Chas's rating
      of Titanic with your answers to Question~2(b).
      \begin{answer}
\begin{alltt}
R = np.array([[2, np.nan, 4, 5, 2],\\
\hspace*{7em}[5, 3, np.nan, 3, np.nan],\\
\hspace*{7em}[np.nan, 2, 3, 5, 3]])\\
W = np.logical\_not(np.isnan(R))\\
K, eta = 1, 0.1\\
rng = np.random.default\_rng(0)\\
A, B = rng.random((3, K)) + 1, rng.random((K, 5)) + 1\\
for it in range(200):\\
\hspace*{2em}for i in range(3):\hspace*{2em}\# ridge for row i of A\\
\hspace*{4em}Bi, Ri = B[:, W[i]], R[i, W[i]]\\
\hspace*{4em}G = eta * np.eye(K) + Bi @ Bi.T\\
\hspace*{4em}A[i] = np.linalg.solve(G, Bi @ Ri)\\
\hspace*{2em}for j in range(5):\hspace*{2em}\# ridge for column j of B\\
\hspace*{4em}Aj, Rj = A[W[:, j]], R[W[:, j], j]\\
\hspace*{4em}G = eta * np.eye(K) + Aj.T @ Aj\\
\hspace*{4em}B[:, j] = np.linalg.solve(G, Aj.T @ Rj)\\
print(np.round(A @ B, 2))
\end{alltt}
        This predicts $3.52$ for Bob and The Matrix and $3.66$ for Chas
        and Titanic, against $3.80$ and $3.39$ from the means.  Other
        random starts give the same answer.  With only 11 ratings and
        $3+5=8$ parameters there is little to learn: without subtracting
        the means, a rank-1 model mostly learns how highly each film is
        rated overall (Spartacus highest, Dr No and Ten lowest), and the
        training RMS error is still $0.88$ because Alice and Bob disagree
        about Titanic.  Matrix factorisation pays off on large, sparse
        data sets.
      \end{answer}}{\lines{9}}

  \end{qparts}
\end{question}

\end{document}
